
\subsection*{2.4}
\phantomsection
\label{sga3:VIIB:2.4}
\addcontentsline{toc}{subsection}{2.4. Quotients of formal groups}

We now return to an arbitrary pseudocompact ring \(k\), in order to
apply to formal groups the results of \SGARefLink{sga3:VIIB:1.4.thm}{1.4}--\SGARefLink{sga3:VIIB:1.5}{1.5} on passage to the
quotient by a topologically flat equivalence relation.
\footnote{\label{sga3:VIIB:nde:90}The editors supplied the details of
the original argument that follow.}

Let \(u:H\to G\) be a monomorphism of formal \(k\)-groups, let
\(\mu:G\times G\to G\) be the multiplication morphism of \(G\), and
let \(\lambda\) be the composite
\begin{center}
  \phantomsection
  \label{sga3:VIIB:2.4:lambda-composite}
  \(
\lambda:\quad
G\times H\xrightarrow{\ \operatorname{id}_G\times u\ }
G\times G\xrightarrow{\ \mu\ }G.
\)

\end{center}
Since \(u\) is a monomorphism, the pair
\begin{center}
  \phantomsection
  \label{sga3:VIIB:2.4:equivalence-pair}
  \begin{tikzpicture}[baseline=(lhs.base)]
  \node (lhs) {$G\times H$};
  \node[right=2.5cm of lhs] (rhs) {$G$};
  \draw[-{Straight Barb[length=4pt,width=4pt]}]
    ([yshift=2.5pt]lhs.east) --
    node[above] {$\operatorname{pr}_1$}
    ([yshift=2.5pt]rhs.west);
  \draw[-{Straight Barb[length=4pt,width=4pt]}]
    ([yshift=-2.5pt]lhs.east) --
    node[below] {$\lambda$}
    ([yshift=-2.5pt]rhs.west);
\end{tikzpicture}

\end{center}
is an equivalence pair in \(\operatorname{Vaf}/k\)
\textup{(cf.\ V, 2.b)}. Recall \textup{(cf.\ \SGARefLink{sga3:VIIB:1.2.C}{1.2.C})} that the
cokernel \(G/H\) of this pair is defined as follows.

Let \(\mathcal O(G)\) and \(\mathcal O(H)\) be the affine algebras of
\(G\) and \(H\), let
\[
  \Delta:\mathcal O(G)
  \longrightarrow
  \mathcal O(G)\widehat\otimes_k\mathcal O(G)
\]
be the diagonal morphism of \(\mathcal O(G)\), and let \(I\) be the
kernel of
\[
  u^\natural:\mathcal O(G)\longrightarrow\mathcal O(H).
\]
By \SGARefLink{sga3:VIIB:1.3.prop}{Proposition 1.3}, the morphism \(u^\natural\) induces an
isomorphism
\[
  \mathcal O(G)/I\xrightarrow{\ \sim\ }\mathcal O(H).
\]
The affine algebra \(\mathcal O(G/H)\) of \(G/H\) is then the kernel
of the pair
\begin{center}
  \phantomsection
  \label{sga3:VIIB:2.4:affine-equalizer-pair}
  \begin{tikzpicture}[baseline=(lhs.base)]
  \node (lhs) {$\mathcal O(G)$};
  \node[right=3cm of lhs] (rhs)
    {$\mathcal O(G)\widehat\otimes_k\mathcal O(H),$};
  \draw[-{Straight Barb[length=4pt,width=4pt]}]
    ([yshift=2.5pt]lhs.east) --
    node[above] {$\tau_1$}
    ([yshift=2.5pt]rhs.west);
  \draw[-{Straight Barb[length=4pt,width=4pt]}]
    ([yshift=-2.5pt]lhs.east) --
    node[below] {$(\operatorname{id}\widehat\otimes u^\natural)\Delta$}
    ([yshift=-2.5pt]rhs.west);
\end{tikzpicture}

\end{center}
where \(\tau_1(x)=x\widehat\otimes1\). Thus
\[
  \mathcal O(G/H)
  =
  \bigl\{
    x\in\mathcal O(G)
    \mathrel{\big|}
    \Delta(x)-x\widehat\otimes1
    \in\mathcal O(G)\widehat\otimes I
  \bigr\}.
\]

If, moreover, \(H\) is topologically flat over \(k\), then
\(\operatorname{pr}_1\) is topologically flat, and \SGARefLink{sga3:VIIB:1.4.thm}{Theorem 1.4} gives
the following result.

\medskip
\noindent\textbf{Theorem.}
\phantomsection
\label{sga3:VIIB:2.4:theorem}
{\itshape Let \(u:H\to G\) be a monomorphism of formal
\(k\)-groups. Suppose that \(H\) is topologically flat over \(k\).
Then the projection
\[
  p:G\longrightarrow G/H
\]
is surjective and topologically flat, there is an isomorphism
\[
  G\times H
  \xrightarrow{\ \sim\ }
  G\times_{G/H}G,
  \tag{*}\label{sga3:VIIB:2.4:equation:quotient-fiber-product-star}
\]
and \(G/H\) represents the quotient sheaf for the flat topology.

Consequently, \(G/H\) has a canonical structure of an object on which
\(G\) acts, and \(p:G\to G/H\) is an equivariant morphism. If, in
addition, \(u\) identifies \(H\) with a normal subgroup of \(G\),
then \(G/H\) has a canonical structure of formal \(k\)-group for
which \(p:G\to G/H\) is a morphism of formal \(k\)-groups, and \(H\)
is the kernel of \(p\).}

Indeed, the first assertion follows from \SGARefLink{sga3:VIIB:1.4.thm}{1.4}, and the other two from
\SGARefLink{sga3:iv:corollary:5:2:2}{Exposé IV, Corollaries 5.2.2} and \SGARefLink{sga3:VIIB:5.2.4}{5.2.4}.

\medskip
\noindent\textbf{Corollary.}
\phantomsection
\label{sga3:VIIB:2.4:corollary}
\footnote{\label{sga3:VIIB:nde:91}The editors made this corollary
explicit; it is used, for example, in \SGARefLink{sga3:VIIB:5.2.1:section-proposition}{5.2.1} and \SGARefLink{sga3:VIIB:5.2.3}{5.2.3}.}
{\itshape Let \(G\) be a formal \(k\)-group and \(H\) a formal
subgroup of \(G\). Let \(A\), \(A/J\), and \(B\) be the affine
algebras of \(G\), \(H\), and \(G/H\), respectively; let \(I_A\) be
the augmentation ideal of \(A\), and put \(I_B=B\cap I_A\). Suppose
that \(H\) is topologically flat over \(k\). Then
\[
  J=AI_B,
\]
the closed ideal generated by \(I_B\).}

Indeed, the projection \(B\to B/I_B\) corresponds to the unit section
\[
  e:\operatorname{Spf}(k)\longrightarrow G/H
\]
of \(G/H\). By \textup{(*)}, \(H\) is identified with the fiber
product
\[
  \operatorname{Spf}(k)\times_{G/H}G.
\]
Its affine algebra \(A/J\) is therefore identified with
\[
  (B/I_B)\widehat\otimes_BA
  \simeq
  A/AI_B.
\]
